\section{Results}
This section shows the findings derived from the development iterations and the overall product benchmarks defined in the previous section.
The derived data is organized to directly address each research question before finally relating the findings to the hypothesis.

\subsection{RQ1: Genetic Algorithm Performance}
To answer RQ1 and whether a GA can be used to create dungeon levels with good flow, we analyzed the evolutionary convergence and the solvability
of the generated levels. By using its Fitness function to optimize its population over time, the GA was able to consistently converge on playable levels.
Initial populations would always be largely unplayable with a low fitness score, but after roughly 30 generations it was able to output playable levels.
Further generations saw continous improvements in terms of fitness before reaching its maximum fitness after 100 generations. Due to the fitness function,
the GA managed to achieve 100\% solvability confirming its ability to filter out impossible layouts. The GA generators difficulty accuracy was not
perfect, but it was good enough to be deemed a reliable offline DDA system.

\subsection{RQ2: Comparative Analysis of GA vs RL}
To answer RQ2 and how GAs and RL agents compare in a PCG context, we compared the two generators across 3 metrics in a benchmark test.
\begin{enumerate}
  \item \textbf{Speed:} The largest difference between the generators were their speeds. The GA required 12 seconds on average to generate a single valid level. 
  This is explained by the iterative and evolutionary search performed by the algorithm, as it needs to simulate gameplay on optimize over a large population size
  and number of generations.
  Compared to the RL agent, it was far behind, as the RL agent relied on inference from its training and managed to generate levels in just 0.125 seconds. 
  \item \textbf{Solvability:} This is where the GA proved superior as it achieved 100\% solvability compared to the RL agent only achieving 91.2\% solvability.
  \item \textbf{Difficulty Accuracy:} The difficulty accuracy is the absolute difference between a levels intended difficulty and the deterministic player 
  agents actual reward. While this is not a perfect solution to determining the "flow" of a level, it proved as a reliable baseline metric to make a difficulty assesment.
  Here the RL agent pulled ahead of the GA with a reward delta of 8.7 vs 13.5. However, this value is slightly biased as the average was only calculated from
  the solvable levels, meaning the actual reward delta for the RL agent is worse over 1000 generated levels. 
\end{enumerate}

\subsection{Hypothesis Findings}
Regarding the initial hypothesis of a GA being better than stochastic methods at creating engaging dungeon layouts, our findings support this fact.
By providing 100\% solvability and a reasonable difficulty accuracy, the GA proves to be much better than stochastically generated levels akin to those
seen in early GA populations. However, due to the computational requirements of the algorithm, it might be more suitable to use a RL approach if runtime
generation is a requirement for the product usecase.

\subsection{Data Summary}
The results of the benchmarks are presented in Table \ref{tab:benchmark-results} below.

\begin{table}[htbp]
  \caption{Benchmark Results}
  \begin{center}
    \begin{tabular}{|c|c|c|}
      \hline
      \textbf{Metric}                           & \textbf{GA} & \textbf{RL}   \\
      \hline
      \textbf{Avg. Generation Speed}            & 12 seconds  & 0.125 seconds \\
      \hline
      \textbf{Solvability Rate}                 & 100\%       & 91.2\%        \\
      \hline
      \textbf{Avg. Reward Delta}$^{\mathrm{a}}$ & 13.5        & 8.7           \\
      \hline
      \multicolumn{3}{l}{$^{\mathrm{a}}$ $|Difficulty Target - Observed Reward|$. Lower is better.}
    \end{tabular}
    \label{tab:benchmark-results}
  \end{center}
\end{table}

% \begin{itemize}
%   \item \textbf{Speed Comparison:} The RL generator achieved an average generation time of \textbf{0.125 seconds} per level and the GA generator
%         achieved
%         This represents a 96\% speed increase over the GA baseline of 12 seconds, validating the main motivation for this iteration.
%   \item \textbf{Solvability:} The trained agent achieved a solvability rate of 91.2\%. This is lower than the GA version
%         that guarantees solvability with its fitness function, although it is sufficient for a system that can regenerate invalid levels in milliseconds.
%         Additionally, the agent achieved an average delta reward of 8.699, meaning the target difficulty wasn't always reached by the generator.
% \end{itemize}

% More thorough results chapter than the one from methods/implementations

% Show statistics if any and respond to the Research questions and/or hypothesis.

% Is the problem resolved?
% How have you resolved the problem/s
% Show results preferably some statistic or other form of numbers.
